% AI-written, AI cross-reviewed; no human sign-off; not a claim that RH is proved
% RH Notes blueprint source excerpts for Z42-SITES
% Verbatim excerpts in the notation of the parent blueprint. The labelled editions and source ranges appear below.
% See the unit YAML for the current proof status. Earlier-round passages retain their original wording.
\section*{Z42-SITES. Served-site parameters and the lattice dictionary}

% Parent source SHA-256: 1af64922b98ed2310744411b1f18b25cc426c97220646321471836e1a3a409d6
% blueprint/ch_P.tex lines 12384 through 12397; r5:Z4.2.01
\subsection{Z4.2.01. The parameters at a served site}\label{r5:Z4.2.01}
At $(D,u,w,h_g)$ set $p=\operatorname{resChar}(w)$, $k_1=P.m$ and write $m_f$ for the fill length $\mathrm{Fl.mm}$. Fix a position $v$. We verify the eleven fields of $\mathrm{Std}$ for the actual local model at $v$.
From $h_g$ we have primality and $5\le p$. The local setting gives even $n'_v\ge2$. The lattice tie gives $m_f=D.R.m$ and the formula for $r$ below. The rank datum gives $D.R.m\ge2$, and $\mathrm{GatTies}.\mathrm{gat\_k}$ gives $k_1+1=D.R.k$. Thus
\[
 r=\begin{cases}r_G& p<p_0(D.R.k),\\1& p\ge p_0(D.R.k).\end{cases}
 \qquad 1\le r_G=K_G+1,\qquad k_1+1=D.R.k.
\]
Both branches satisfy $1\le r$ and $2r\le k_1+1$ by the explicit hypothesis $2r_G\le D.R.k$. This also proves the required $m_f\ge1$. Take the base-weight identity from the fill and dominance from the selected local model. The remaining weight fields require $c\le w_1(s)$ and $-\lambda_0(s,k_1)\le w_1(s)$, both consequences of the fill's single bound for their maximum. Its monotonicity holds for every $j\ge1$ and hence for all $1\le j<m_f$. This uses the base weight, terminal coordinate and sequence from the same position.

{\footnotesize \nolinkurl{GlobalTheory.TwU.std_at_served} (\nolinkurl{Development/AbelianVarietyTheory/Proofs/SatakeParameters/StdAtServed.lean}\allowbreak:9--22). \nolinkurl{Formalism.MainNotions.FormalGroupSystem.Std} (\nolinkurl{Formalism/SharpRestriction/GrowthBounds/FormalGroupSystem.lean}\allowbreak:89--101). \nolinkurl{Formalism.MainNotions.G.rG}, \nolinkurl{LatTied} (\nolinkurl{Formalism/SharpRestriction/GrowthBounds/IsResidueClassZero.lean}\allowbreak:41,430--436). \nolinkurl{Formalism.BasicNotions.GatTies} (\nolinkurl{Formalism/LocalRingCardinality/DataGatTies.lean}\allowbreak:7--10). \nolinkurl{Fills.base}, \nolinkurl{Fills.w_ge}, \nolinkurl{Fills.w_mono} (\nolinkurl{Formalism/SharpRestriction/GrowthBounds/DataGChain.lean}\allowbreak:27--38).}

\emph{Complete for a tied served site.} [Hypothesis, discharged at the application:] the joint object supplies its lattice tie, and the chain theorem receives the rank inequality as its explicit hypothesis. It uses this very $\mathrm{Std}$ to construct the comparison at that site.

{\footnotesize \nolinkurl{GlobalTheory.TwV.site_cmpLite_siteAdm} (\nolinkurl{Development/AbelianVarietyTheory/Proofs/SatakeParameters/JointSiteChainComparedCarriesAdmittedSystem.lean}\allowbreak:77--89).}

% blueprint/ch_P.tex lines 12533 through 12548; r5:small:Z4.2.03
\subsection{Z4.2.03. One adelic point for the class, kernel lattice and determinant}
\label{r5:small:Z4.2.03}
\textit{Complete under the given standard joint.} Fix the tagged datum $t$, the comparison hypothesis and labelling, the standard local model and the joint $J$. For every $j\le m$ and valid tuple $L$, let $x=J.\mathrm{pt}(j,L)$. The four statements concern this same $x$: its defining point condition; $x=\mathrm{rat}(\gamma)y_{c(L)}\kappa$ for a rational unitary $\gamma$ and a member $\kappa$ of the level group of the selected chart; identification after rank transport of the kernel lattice with $\mathrm{latK}(\mathrm{gram}_j,J.a,x,L_b)$; and equality of the frame determinant with $\det x$.
\par{\footnotesize \nolinkurl{GlobalTheory.TwV.joint_dict} (\nolinkurl{Development/AbelianVarietyTheory/Proofs/SatakeParameters/LatticeTupleJointDictionary.lean}\allowbreak:16--25). \nolinkurl{GlobalTheory.TwV.Joint.Dict} (\nolinkurl{Development/AbelianVarietyTheory/Definitions/GaloisGeometry/JointCmpEmbeddingY.lean}\allowbreak:123--129).\par}

At each place the tuple condition exhibits a matrix of invertible determinant spanning the local lattice. It is a member of $\mathrm{GL}$ and witnesses the defining choice of local bases. A good chart has a section for the prescribed local matrices. The product of its fixed unitary representative and the section is a unitary point with precisely these local frames. In level zero use the good chart from $J.\mathrm{one}$. At $j=i+1$ split the component via the inverse of the component equivalence into its chart $\beta$ and its local component, and use the goodness of $H.\mathrm{core}(i,\beta)$. In either case the chart construction already chose such a point by classical choice. To identify the point chosen in $J$ with the chart point, unfold the defining point condition: it specifies the chart matrix after a rank reindexing. Applying the inverse reindexing gives back that very matrix entry by entry. Thus its local conditions hold for $x$.
\par{\footnotesize \nolinkurl{GlobalTheory.TwV.Joint.basesAt_span} (\nolinkurl{Development/AbelianVarietyTheory/Proofs/SatakeParameters/LatticeTupleBasesSpanAtPlaces.lean}\allowbreak:10--16). \nolinkurl{GlobalTheory.TwU.exists_unitary_pt} (\nolinkurl{Development/AbelianVarietyTheory/Proofs/SatakeParameters/ExistsUnitaryPt.lean}\allowbreak:14--16). \nolinkurl{GlobalTheory.TwV.Joint.isPt_pt_zero} (\nolinkurl{Development/AbelianVarietyTheory/Proofs/SatakeParameters/LatticeTupleAdelicPointStepZero.lean}\allowbreak:12--35). \nolinkurl{GlobalTheory.TwV.Joint.pt_succ_isPt} (\nolinkurl{Development/AbelianVarietyTheory/Proofs/SatakeParameters/LatticeTupleAdelicPointStepSucc.lean}\allowbreak:14--41).\par}

By the surjectivity of chart representatives onto the double quotient, one such local point lies in $\mathrm{rat}(\gamma)y_c U$. For two choices of bases spanning the same lattice, the ratio of their points is the identity off $p$. Over $p$ their change-of-basis matrix and its inverse preserve the standard lattice. After conjugation by the prescribed frame, the chart hypotheses \path{LocInt} and \path{LocHyp}, with the bad-prime condition, put the ratio in the local level group. Thus their adelic ratio lies in $U$. This proves that the chosen $x$ works with the class and rational element selected from those bases. At zero $J.\mathrm{cls0}$ identifies the class with the image under $J.e_0$ of the level-zero class; $J.y_0$ says its displayed representative differs from the chart representative by some $\kappa'\in U$ on the right; replace the right factor $\kappa$ by $(\kappa')^{-1}\kappa$. At a successor the chart of $c(L)$ is $\beta$. The component equivalence and its inverse cancel, and the chart transport identifies the two level groups and their representatives. This puts exactly $y_{c(L)}$ in the asserted equation.
\par{\footnotesize \nolinkurl{Formalism.L2.G.ptEx_of} (\nolinkurl{Formalism/LocalToGlobal/GroupLattices/TiedFamilies/EqOneOfPh.lean}\allowbreak:120--132). \nolinkurl{Formalism.L2.G.ptOne_of} (\nolinkurl{Formalism/LocalToGlobal/GroupLattices/TiedFamilies/TiePtOne.lean}\allowbreak:55--83). \nolinkurl{Formalism.MainNotions.G.KAat.IsCG1.pt} (\nolinkurl{Formalism/SharpRestriction/GrowthBounds/ChartedClasses/LocalClassDataOfTrivialStabilizers.lean}\allowbreak:90--101). \nolinkurl{Formalism.MainNotions.G.KAat.kaAt_pt} (\nolinkurl{Formalism/SharpRestriction/GrowthBounds/ChartedClasses/ClsChartedLabAt.lean}\allowbreak:141--157). \nolinkurl{GlobalTheory.TwV.Joint.pt_zero_eq_rat_mul_rep_mul} (\nolinkurl{Development/AbelianVarietyTheory/Proofs/SatakeParameters/StepZeroPointMemDoubleCosetClassRep.lean}\allowbreak:12--26). \nolinkurl{GlobalTheory.TwV.Joint.pt_succ_coset} (\nolinkurl{Development/AbelianVarietyTheory/Proofs/SatakeParameters/SuccStepPointMemDoubleCosetClassRep.lean}\allowbreak:13--30).\par}

For the kernel lattice at zero, $J.\mathrm{lev0}$ identifies the base lattice. At a successor, unfold the chart for $i+1\le m$: its Gram form and its base lattice are the transported form and lattice of the component chart. Its frame is the reindexed matrix already identified with $x$. Transport all three objects through the same dimension equality. Substituting equal ranks reduces matrix reindexing and the transport of lattice vectors to their identities, so taking the kernel lattice before or after transport gives the same submodule. The determinant of the frame is by definition the determinant of the chart matrix. Simultaneous row and column reindexing preserves this determinant, proving the last clause.
\par{\footnotesize \nolinkurl{GlobalTheory.kerN_zero_eq_latK_pt} (\nolinkurl{Development/AbelianVarietyTheory/Proofs/SatakeParameters/StepZeroKernelLatticeEqPointKernelLattice.lean}\allowbreak:12--25). \nolinkurl{GlobalTheory.TwV.Joint.kerN_succ_heq} (\nolinkurl{Development/AbelianVarietyTheory/Proofs/SatakeParameters/SuccStepKernelLatticeEqPointKernelLattice.lean}\allowbreak:63--72). \nolinkurl{GlobalTheory.TwV.chart_succ_eq} (\nolinkurl{Development/AbelianVarietyTheory/Proofs/SatakeParameters/LatticeTupleChartSuccEqComponentChart.lean}\allowbreak:22--27). \nolinkurl{GlobalTheory.reindex_rename_point_det_latK} (\nolinkurl{Development/AbelianVarietyTheory/Proofs/SatakeParameters/ReindexRenamePointDetLatK.lean}\allowbreak:14--28). \nolinkurl{GlobalTheory.joint_dt_eq_det_of_isPt} (\nolinkurl{Development/AbelianVarietyTheory/Proofs/SatakeParameters/LatticeTupleDetEqAdelicPointDet.lean}\allowbreak:13--18).\par}

\textit{[Hypothesis, discharged at the use:]} the joint and chart hypotheses are the supplied arguments. The site construction obtains standardness from the served-site bound and the tie of that joint, then uses this dictionary in the same admitted-system construction.
\par{\footnotesize \nolinkurl{GlobalTheory.TwV.site_cmpLite_siteAdm} (\nolinkurl{Development/AbelianVarietyTheory/Proofs/SatakeParameters/JointSiteChainComparedCarriesAdmittedSystem.lean}\allowbreak:77--83). \nolinkurl{admissible_at_each_site} (\nolinkurl{Development/AbelianVarietyTheory/Proofs/SatakeParameters/JointSiteChainComparedCarriesAdmittedSystem.lean}\allowbreak:29--34).\par}

% blueprint/ch_P.tex lines 12550 through 12564; r5:small:Z4.2.05
\subsection{Z4.2.05. Descent of position embeddings through conjugation}
\label{r5:small:Z4.2.05}
\textit{Complete for every site.} Fix a position $(v,s)$, write $\tau=\mathrm{pos}(v,s)$, and put $x_s=x$ for even $s$, $x_s=\bar x$ for odd $s$. The assertion $S.\mathrm{Pos}$ is the identity
\[
  \mathrm{lamC}_{\lambda}\bigl(\mathrm{embT}_{\tau}(x)\bigr)
     =\mathrm{twG}_{v,s}\bigl(\rho_v(x_s)\bigr).
\]
The two source hypotheses are the compatibility $i_K(ez)=\lambda(\mathrm{embLL}(z))$ on algebraic arguments, and the position identity for $e(\mathrm{tEL}_{\tau}(y))$ with its even/odd involution on the input. Both use the same local field embedding $e$.
\par{\footnotesize \nolinkurl{GlobalTheory.TwV.pos_at_site} (\nolinkurl{Development/AbelianVarietyTheory/Proofs/SatakeParameters/SitePlaceInducesPositionEmbeddings.lean}\allowbreak:11--13). \nolinkurl{GlobalTheory.TwV.Joint.At} (\nolinkurl{Development/AbelianVarietyTheory/Definitions/GaloisGeometry/JointCmpEmbeddingY.lean}\allowbreak:134--138). \nolinkurl{GlobalTheory.TwV.Joint.At.Pos} (\nolinkurl{Development/AbelianVarietyTheory/Definitions/GaloisGeometry/JointCmpEmbeddingY.lean}\allowbreak:157--160). \nolinkurl{GlobalTheory.TwU.LamPos} (\nolinkurl{Development/AbelianVarietyTheory/Definitions/GaloisGeometry/NatCastMemOfPlace.lean}\allowbreak:53--54).\par}

Apply the position identity to $x_s$. At even $s$ its argument on the left is $x$. At odd $s$ its argument is $\mathrm{sigD}(\bar x)=x$, since $\mathrm{sigD}$ is the host conjugation and its square is the identity. The result in the local field is $e(\mathrm{tEL}_{\tau}(x))=\iota_v\sigma_v^s\rho_v(x_s)$. The element $z=\mathrm{tEL}_{\tau}(x)$ belongs to the number field $LL$, so its complex image is algebraic. By compatibility of the two global embeddings this complex image is $\mathrm{embT}_{\tau}(x)$. On an algebraic argument $\mathrm{lamC}$ evaluates the given $\lambda$ at that very element of the algebraic closure. Use the site compatibility to identify it with $i_K(ez)$ and substitute the position identity. This gives the displayed equality for all $x,v,s$.
\par{\footnotesize \nolinkurl{Formalism.L2.G.embLL_mem_alg} (\nolinkurl{Formalism/LocalToGlobal/GroupLattices/LevelGatherings/PosBridge.lean}\allowbreak:24--26). \nolinkurl{Formalism.L2.G.Cmp.tEL_emb} (\nolinkurl{Formalism/LocalToGlobal/GroupLattices/ComparisonMaps/TiedRec.lean}\allowbreak:18--26). \nolinkurl{GlobalTheory.TwU.PlaceAbove.involution_conj} (\nolinkurl{Development/AbelianVarietyTheory/Proofs/SatakeParameters/LamPosOfPosEmb.lean}\allowbreak:20--22). \nolinkurl{GlobalTheory.TwU.PlaceAbove.posArg_eq_self} (\nolinkurl{Development/AbelianVarietyTheory/Proofs/SatakeParameters/LamPosOfPosEmb.lean}\allowbreak:24--33). \nolinkurl{GlobalTheory.TwU.lamPos_of_posEmb} (\nolinkurl{Development/AbelianVarietyTheory/Proofs/SatakeParameters/LamPosOfPosEmb.lean}\allowbreak:39--54).\par}

\textit{[Hypothesis, discharged at the use:]} the site supplies $S.hK$ and the chosen joint supplies the position identity through its comparison tie. The admitted-system construction applies this equality to its very same site before invoking the arithmetic frame transformation.
\par{\footnotesize \nolinkurl{GlobalTheory.TwV.pos_at_site} (\nolinkurl{Development/AbelianVarietyTheory/Proofs/SatakeParameters/SitePlaceInducesPositionEmbeddings.lean}\allowbreak:11--13). \nolinkurl{admissible_at_each_site} (\nolinkurl{Development/AbelianVarietyTheory/Proofs/SatakeParameters/JointSiteChainComparedCarriesAdmittedSystem.lean}\allowbreak:29--34).\par}

% blueprint/ch_P.tex lines 12615 through 12617; r6:review:Z4.2.01
\paragraph{Z4.2.01. Served-site parameters (complete).}\phantomsection\label{r6:review:Z4.2.01} \emph{Original argument on p.~\pageref{r5:Z4.2.01}.} 
The rank data give $2\le m$, the lattice tie identifies $m_f=m$ and specifies the two values of $r$, and the separate gate tie gives $k_1+1=k$. In the $r=1$ branch, $1\le r_G$ and the explicit assumption $2r_G\le k$ supply $2\le k$. The caller supplies this inequality and the selected joint's lattice tie. The local setting, fill and served-site hypotheses provide the other fields as in Section~\ref{r5:Z4.2.01}.
See \path{std_at_served}, \path{site_cmpLite_siteAdm} (\path{Development/AbelianVarietyTheory/Proofs/SatakeParameters/StdAtServed.lean:9--22}, \path{Development/AbelianVarietyTheory/Proofs/SatakeParameters/JointSiteChainComparedCarriesAdmittedSystem.lean:77--89}).

% blueprint/ch_P.tex lines 12619 through 12621; r6:review:Z4.2.03
\paragraph{Z4.2.03. Four assertions for a fixed joint (complete).}\phantomsection\label{r6:review:Z4.2.03} \emph{Original argument on p.~\pageref{r5:small:Z4.2.03}.} 
For each lattice tuple, the same selected point satisfies the local point condition, represents the rational-unitary/level double coset, has the specified kernel lattice after rank transport and has the specified determinant. At zero, $J.\mathrm{cls0}$ identifies the class label, $J.y_0$ compares representatives and $J.\mathrm{lev0}$ identifies the level for the correction $\kappa'^{-1}\kappa$. The right ratio lies in the level group by local integrality, the local level implication and the hypothesis on bad primes. These are assumptions on the given chart and joint.
See \path{Joint.Dict}, \path{joint_dict}, \path{pt_zero_eq_rat_mul_rep_mul}, \path{ptOne_of} (\path{Development/AbelianVarietyTheory/Definitions/GaloisGeometry/JointCmpEmbeddingY.lean:123--129}, \path{Development/AbelianVarietyTheory/Proofs/SatakeParameters/LatticeTupleJointDictionary.lean:16--25}, \path{Development/AbelianVarietyTheory/Proofs/SatakeParameters/StepZeroPointMemDoubleCosetClassRep.lean:12--26}, \path{Formalism/LocalToGlobal/GroupLattices/TiedFamilies/TiePtOne.lean:55--82}).

% blueprint/ch_P.tex lines 12623 through 12624; r6:review:Z4.2.05
\paragraph{Z4.2.05. Position and conjugation (complete).}\phantomsection\label{r6:review:Z4.2.05} \emph{Original argument on p.~\pageref{r5:small:Z4.2.05}.} 
The site supplies $S.hK$ and the comparison tie supplies its position embedding. For odd $s$, insert $\bar x$ before applying the position identity, so that the host involution sends it to $x$. At even $s$ insert $x$. Algebraicity of the image of the resulting number-field element allows the use of the specified embedding $\lambda$. Section~\ref{r5:small:Z4.2.05} proves the identity for every site. See \path{Joint.At.Pos}, \path{pos_at_site}, \path{lamPos_of_posEmb} (\path{Development/AbelianVarietyTheory/Definitions/GaloisGeometry/JointCmpEmbeddingY.lean:134--145,159}, \path{Development/AbelianVarietyTheory/Proofs/SatakeParameters/SitePlaceInducesPositionEmbeddings.lean:11--13}, \path{Development/AbelianVarietyTheory/Proofs/SatakeParameters/LamPosOfPosEmb.lean:20--32,39--54}).

